\documentclass[11pt, a4paper]{article}

\usepackage[a4paper, margin=2.5cm]{geometry}
\usepackage[utf8]{inputenc}
\usepackage[T1]{fontenc}
\usepackage[english]{babel}
\usepackage{lmodern}

\usepackage{amsmath}
\usepackage{amssymb}
\usepackage{booktabs}
\usepackage{listings}
\usepackage{xcolor}
\usepackage{graphicx}
\usepackage{hyperref}

% Academic, restrained styling for listings
\definecolor{codebg}{gray}{0.96}
\definecolor{codeframe}{gray}{0.75}
\definecolor{codegreen}{rgb}{0.1, 0.45, 0.1}
\definecolor{codeblue}{rgb}{0.0, 0.15, 0.55}
\definecolor{codegray}{rgb}{0.35, 0.35, 0.35}

\lstset{
  backgroundcolor=\color{codebg},
  basicstyle=\ttfamily\footnotesize,
  breaklines=true,
  commentstyle=\color{codegreen}\itshape,
  frame=lines,
  rulecolor=\color{codeframe},
  keywordstyle=\color{codeblue}\bfseries,
  stringstyle=\color{codegray},
  numbers=none,
  tabsize=2,
  showstringspaces=false,
  aboveskip=0.8em,
  belowskip=0.8em
}

\hypersetup{
  colorlinks=true,
  linkcolor=black,
  urlcolor=codeblue,
  citecolor=black
}

\begin{document}

\begin{center}
  {\scshape\large Universitat Politècnica de València}\\[0.2em]
  {\scshape\normalsize Escola Tècnica Superior d'Enginyeria Informàtica (ETSINF)}\\[1.2em]
  {\Large\bfseries Practical 1: Introduction to R}\\[0.4em]
  {\large Data Science}\\[1.2em]
  \begin{tabular}{rl}
    \textbf{Author(s):} & Enrique Sopeña Urbano \\
    \textbf{Subject:}   & Practical Work 1      \\
    \textbf{Date:}      & \today
  \end{tabular}
\end{center}

\vspace{0.8em}
\hrule
\vspace{1.2em}

\begin{abstract}
  \noindent This laboratory report documents the exercises completed for Practical~1 of the Data Science course. The objective is to establish familiarity with the R scientific programming environment, covering primitive vector manipulation, cyclic repetitions, matrix construction, pseudorandom generation under controlled seed states, data frame ingestion, frequency cross-tabulations, conditional partitioning, and graphical bivariate inspection.
\end{abstract}

\vspace{1em}

\section{Vector Sequence Generation}
\textbf{Objective:} Generate the sequence of natural numbers $\{1, 2, \dots, 12\}$ and store the resulting array in the vector object $X$.

\subsection{Implementation and Output}
\begin{lstlisting}[language=R]
# Generate the sequence 1 through 12 using the colon operator
X <- 1:12
X
\end{lstlisting}

\begin{lstlisting}
 [1]  1  2  3  4  5  6  7  8  9 10 11 12
\end{lstlisting}

\subsection{Discussion}
In the R language, the binary colon operator (\texttt{:}) acts as a compact generator for regular arithmetic sequences with unit stride. The assignment operator \texttt{<-} allocates the resulting 12-element integer vector to symbol $X$.

\section{Cyclic Sequence Repetitions}
\textbf{Objective:} Construct four contiguous repetitions of the numeric sequence $\mathbf{s} = (6, 2, 4)$.

\subsection{Implementation and Output}
\begin{lstlisting}[language=R]
# Repeat the entire tuple (6, 2, 4) four times
rep(c(6, 2, 4), times = 4)
\end{lstlisting}

\begin{lstlisting}
 [1] 6 2 4 6 2 4 6 2 4 6 2 4
\end{lstlisting}

\subsection{Discussion}
The primitive function \texttt{rep()} replicates input vector components. By setting \texttt{times = 4}, the complete ordered sequence of length 3 is repeated cyclically, producing an atomic numeric vector of length 12 ($3 \times 4$).

\section{Matrix Construction from Sequences}
\textbf{Objective:} Generate a sequence comprising six 9s, five 2s, and four 5s, populating a $5 \times 3$ matrix column-wise.

\subsection{Implementation and Output}
\begin{lstlisting}[language=R]
# Construct concatenated repetition vector
vec3 <- c(rep(9, 6), rep(2, 5), rep(5, 4))

# Organize elements into a 5x3 matrix populated column-wise
M <- matrix(vec3, nrow = 5, ncol = 3, byrow = FALSE)
M
\end{lstlisting}

\begin{lstlisting}
     [,1] [,2] [,3]
[1,]    9    9    2
[2,]    9    2    5
[3,]    9    2    5
[4,]    9    2    5
[5,]    9    2    5
\end{lstlisting}

\subsection{Discussion}
The scalar repetitions form an ordered sequence of cardinality $6 + 5 + 4 = 15$. The \texttt{matrix()} constructor dimensions the array into 5 rows and 3 columns. Because R follows Fortran-style column-major ordering by default (\texttt{byrow = FALSE}), the five rows of column 1 are populated first, followed by columns 2 and 3.

\newpage

\section{Pseudorandom Normal Generation and Reproducibility}
\textbf{Objective:} Generate $N = 20$ independent realizations from a standard normal distribution $\mathcal{N}(0, 1)$ using seed 100, assess central tendency and dispersion metrics, and evaluate seed state invariance.

\subsection{Implementation and Output}
\begin{lstlisting}[language=R]
# Initialize RNG with seed 100 and sample 20 normal values
set.seed(100)
v1 <- rnorm(20)
v1

# Compute descriptive statistics
m_val   <- mean(v1)
med_val <- median(v1)
var_val <- var(v1)
sd_val  <- sd(v1)

c(Mean = m_val, Median = med_val, Variance = var_val, StdDev = sd_val)
\end{lstlisting}

\begin{lstlisting}
> v1
 [1] -0.50219235  0.13153117 -0.07891709  0.88678481  0.11697127  0.31863009
 [7] -0.58179068  0.71453271 -0.82525943 -0.35986213  0.08988646  0.09627446
[13] -0.20163395  0.73984050  0.12337954 -0.02931771 -1.19253457 -0.19819770
[19] -0.08865646  0.64069818

> c(Mean = m_val, Median = med_val, Variance = var_val, StdDev = sd_val)
       Mean      Median    Variance      StdDev 
 0.01499573  0.03028437  0.28004138  0.52918936 
\end{lstlisting}

\subsection{Seed State Comparison}
\begin{lstlisting}[language=R]
# Identical seed evaluation
set.seed(100)
v2 <- rnorm(20)
identical(v1, v2)

# Alternative seed evaluation
set.seed(200)
v3 <- rnorm(20)
c(Mean = mean(v3), StdDev = sd(v3))
\end{lstlisting}

\begin{lstlisting}
> identical(v1, v2)
[1] TRUE

> c(Mean = mean(v3), StdDev = sd(v3))
     Mean    StdDev 
0.0396767 0.9984920 
\end{lstlisting}

\begin{table}[htbp]
  \centering
  \caption{Sample Statistics ($N=20$) versus Theoretical Population $\mathcal{N}(0, 1)$}
  \label{tab:seed_stats}
  \vspace{0.4em}
  \begin{tabular}{lccc}
    \toprule
    \textbf{Statistic} & \textbf{($\text{seed}=100$)} & \textbf{($\text{seed}=200$)} & \textbf{Population Value} \\
    \midrule
    Mean ($\bar{x}$)   & $0.0150$                     & $0.0397$                     & $0.0000$                  \\
    Median             & $0.0303$                     & $-0.0485$                    & $0.0000$                  \\
    Variance ($s^2$)   & $0.2800$                     & $0.9970$                     & $1.0000$                  \\
    Std. Dev. ($s$)    & $0.5292$                     & $0.9985$                     & $1.0000$                  \\
    \bottomrule
  \end{tabular}
\end{table}

\subsection{Discussion}
Pseudorandom number algorithms (such as the default Mersenne Twister in R) generate deterministic sequences conditioned on an internal state register. Specifying \texttt{set.seed(100)} resets the register to a fixed position, yielding strictly identical realization vectors (\texttt{identical(v1, v2) == TRUE}). Modifying the seed to $200$ produces an uncorrelated sequence from the distribution, demonstrating that empirical sample statistics diverge according to finite-sample variation.

\newpage

\section{Data Frame Analysis on Student Observations}
The dataset \texttt{data1.txt} compiles anthropometric measurements and categorical attributes across student subjects.

\subsection{(a--d) Data Ingestion, Header Assignment, and Verification}
\begin{lstlisting}[language=R]
# (a) Read space-delimited text without header
students <- read.table("data1.txt", header = FALSE)

# (b) Assign explicit column titles
colnames(students) <- c("height", "shoesize", "gender", "population")

# (c) Validate structure and inspect leading observations
str(students)
head(students, n = 3)

# (d) Query column titles exclusively
colnames(students)
\end{lstlisting}

\begin{lstlisting}
> colnames(students)
[1] "height"     "shoesize"   "gender"     "population"
\end{lstlisting}

\subsection{(e) Attribute Extraction}
\begin{lstlisting}[language=R]
# Extract single attribute vector
students$height
\end{lstlisting}

\subsection{(f--g) Categorical Distributions and Contingency Analysis}
\begin{lstlisting}[language=R]
# (f) One-way marginal distributions
table(students$gender)
table(students$population)

# (g) Two-way cross-tabulation table
table(students$gender, students$population)
\end{lstlisting}

\begin{table}[htbp]
  \centering
  \caption{Cross-Classification Contingency Matrix (Gender $\times$ Population)}
  \label{tab:contingency}
  \vspace{0.4em}
  \begin{tabular}{lcccc}
    \toprule
    \textbf{Gender} & \textbf{Site 1} & \textbf{Site 2} & \textbf{Site 3} & \textbf{Total} \\
    \midrule
    Female          & $n_{f,1}$       & $n_{f,2}$       & $n_{f,3}$       & $N_f$          \\
    Male            & $n_{m,1}$       & $n_{m,2}$       & $n_{m,3}$       & $N_m$          \\
    \midrule
    \textbf{Total}  & $N_{\cdot,1}$   & $N_{\cdot,2}$   & $N_{\cdot,3}$   & $N$            \\
    \bottomrule
  \end{tabular}
\end{table}

\subsection{(h) Stratification by Gender}
\begin{lstlisting}[language=R]
# Bracket indexing notation
males_df   <- students[students$gender == "male", ]
females_df <- students[students$gender == "female", ]

# Declarative subset function
males_sub   <- subset(students, gender == "male")
females_sub <- subset(students, gender == "female")
\end{lstlisting}

\subsection{(i) Stratification Relative to Median Height}
\begin{lstlisting}[language=R]
# Compute sample median threshold
med_h <- median(students$height)

# Partition using indexing and subsetting
below_df  <- students[students$height < med_h, ]
above_df  <- students[students$height > med_h, ]

below_sub <- subset(students, height < med_h)
above_sub <- subset(students, height > med_h)
\end{lstlisting}
\textit{Remark:} Strict ordering conditions ensure that boundary observations ($x_i = \text{median}$) are excluded from both strict partitions.

\newpage

\subsection{(j) Scale Transformations: Centimetres to Metres}
Three distinct programming paradigms were implemented to apply the metric scale transformation $h_{\text{m}} = h_{\text{cm}} / 100$:

\begin{lstlisting}[language=R]
# Paradigm 1: Vectorized primitive division (idiomatic)
st_m1 <- students
st_m1$height <- st_m1$height / 100

# Paradigm 2: Explicit procedural loop
st_m2 <- students
for (i in 1:nrow(st_m2)) {
  st_m2$height[i] <- st_m2$height[i] / 100
}

# Paradigm 3: Functional programming via apply()
st_m3 <- students
st_m3$height <- apply(st_m3["height"], 1, function(x) x / 100)
\end{lstlisting}

\begin{itemize}
  \item \textbf{Vectorized Arithmetic:} Leverages compiled low-level routines without interpreter overhead, representing the optimal design pattern in R.
  \item \textbf{Procedural Iteration:} Requires repeated row-level indexing in interpreted memory, incurring higher computational complexity $\mathcal{O}(n)$.
  \item \textbf{Functional Application:} Encapsulates transformation per row margin through \texttt{apply()}, guaranteeing consistency across array dimensions.
\end{itemize}

\subsection{(k) Bivariate Scatter Visualization}
\begin{lstlisting}[language=R]
# Encode visual aesthetics conditionally
point_colors <- ifelse(students$gender == "male", "blue", "magenta")
point_shapes <- ifelse(students$gender == "male", 1, 3)

# Render scatter plot
plot(students$shoesize, students$height,
     col = point_colors,
     pch = point_shapes,
     xlab = "Shoe Size",
     ylab = "Height (cm)",
     main = "Height vs Shoe Size by Gender")

# Attach explanatory legend
legend("topleft",
       legend = c("Male", "Female"),
       col = c("blue", "magenta"),
       pch = c(1, 3))
\end{lstlisting}

\vspace{1.5em}
\begin{figure}[htbp]
  \centering
  \includegraphics[width=1\textwidth]{grafico.png}
  \caption{Scatter plot of student shoe size against measured height, stratified by gender.}
  \label{fig:scatter}
\end{figure}

\end{document}